% 第 49 章　相对论性量子力学
\chapter{相对论性量子力学}

在前面几章里，我们先后拿到了两件各自威力无穷的武器：狭义相对论告诉我们，能量、动量和质量被 \(E^2=p^2c^2+m^2c^4\) 牢牢绑在一起（第 40 章）；量子力学告诉我们，微观世界的状态要用波函数描述，它的演化由薛定谔方程支配（第 43、44 章）。两件武器各自在自己的战场上从未失手。一个自然的问题是：把它们缝在一起，会发生什么？本章的答案是：缝合处迸出了火花，火花里掉出了三样没人预订的东西——反物质、自旋，以及一个彻底改写“粒子”含义的新世界观。

\step{反常识开场}

你去医院做一次 PET 检查（正电子发射断层扫描，一种常见的肿瘤与脑功能影像检查），流程听起来很平常：护士往你的静脉里注射一针含氟的葡萄糖类似物，你安静地躺进一个圆环形的机器，十几分钟后拿到一张显示身体各部位代谢活跃程度的彩色图。

但这针药水一点都不平常。其中的氟是一种特殊的同位素氟-18，它的原子核会自发地放出一个“正电子”——一个质量和电子完全相同、电荷符号却相反的粒子。物理学家给它起了个更戏剧性的名字：电子的反物质。正电子在你的组织里穿行不到一毫米，就会撞上某个普通电子。接下来发生的事情违反你从小到大接受的一切常识：两个粒子同时消失了，原地没有留下任何碎屑，只留下两束伽马射线，沿几乎严格相反的方向飞出，每一束带走 511 keV 的能量。机器外面的探测器圆环记下这两束射线到达的精确时刻，计算机把数百万次这样的事件反推回它们的出发位置，就画出了那张彩色图。

也就是说，你的体检报告背后藏着这样一行事实：有反物质在你体内产生，然后与物质相遇，成对湮灭成了光。初中学过“物质不能凭空消失”，化学课反复强调“反应前后原子守恒”——可在这里，一对实实在在的粒子确实消失了，只剩光。更让人坐不住的是这个粒子的来历：正电子不是实验家先在角落里撞见的怪东西，而是 1928 年一个二十七岁的理论家坐在桌前，试图把量子力学和相对论缝进同一个方程时，方程自己“多出来”的解。理论先断言“世界上应该有它”，四年后实验家才在云室照片里真的找到了它。

一个方程怎么会预言一种没人见过的粒子？两个各自完美的理论硬拼在一起，为什么不是得到一个更完美的理论，而是先得到了一堆矛盾，然后才是新大陆？这就是本章要走的路。

\step{先猜一猜}

在往下读之前，请先把你自己的预测写下来。物理学家面对两个成功的理论时，本能的反应和你一样：打个补丁就行了吧。

第一个猜测：薛定谔方程里的动能项是 \(p^2/2m\)，这只是低速近似。那么把这一项换成相对论的正确版本，是不是就得到“相对论版薛定谔方程”，从此万事大吉？你觉得：（a）这样直接换就能成功；（b）会出点小毛病，但修修补补总能凑合；（c）会撞上根本性的墙，怎么补都没用。

第二个猜测：一只被加速到接近光速的电子，我们还能不能像第 43 章那样，把它当成“一个”粒子来描述——给它一个波函数，追踪它的概率分布？换个问法：当碰撞的能量足够高时，“世界里总共有多少个粒子”这件事，是固定不变的花名册，还是会临时改写？

第三个猜测：如果新理论真的预言了某种闻所未闻的粒子，你觉得物理学家会立刻相信方程，还是会先怀疑自己算错了？（想一想你解方程得到一个负数解、而题目问的是人数时，你的第一反应是什么。）

把你的答案记在页边。本章结束时回头看，你会发现：第一题的正确答案接近（c），第二题的答案是“花名册会被改写”，而第三题，连狄拉克本人都犹豫了三年。

\step{动手实验}

本章的现象发生在 511 keV 的能量尺度上，家里不可能直接演示粒子湮灭。但本章最核心的观念翻转——“粒子的数目不是守恒的，场的激发才是基本的”——可以在一盆水里看到它的影子。

\begin{experiment}[一盆水里的“粒子个数”]
材料：一个平底大盘子或水槽，清水，一根筷子，一枚硬币。

第一步，等水面完全平静。用筷子尖在水面快速点一下：你看到一圈整齐的涟漪向外跑——一个“波包”，像一个匆匆赶路的小东西。

第二步，换成用整个手掌拍一下水面。这次出现的不是一圈，而是一大群大大小小、方向各异的涟漪，它们互相穿过、分裂、合并。拍的力度越大，涟漪的数目和花样就越多。

第三步，在水底放一枚硬币制造一块浅水区，再点一下筷子。涟漪经过硬币上方时会绕射、分裂成几股，过了浅水区又重新叠加。

现在数一数：这盆水里“有几个波”？你会发现这个问题没有确定的答案——一个波包可以散成一片，一片波纹也可以聚成一个鼓包。波的“个数”根本不是守恒量，守恒的是你拍进去的能量去了哪里。
\end{experiment}

这个比喻像什么：它像本章结尾要给出的世界观——水面是“场”，涟漪是场的“激发”，激发的数目随你注入的能量多少而改变，没有一个“涟漪总数守恒定律”。它不像什么：水波的能量可以任意细分，而量子场的激发是一份一份的，每一份的大小由普朗克常数钉死；水波衰减成热，量子场的激发却可以真正凭空成对出现、成对消失。一盆水只是借来的直觉，不是证据。

如果条件许可，还有一个更贴近本章主题的家庭项目：用干冰和酒精自制一个简易威尔逊云室（网上有大量教程，注意干冰防冻伤）。在黑暗里用手电照亮云室，你会看到不时划过的白色细线——那是来自太空的宇宙线留下的径迹，其中大部分是 \(\mu\) 子。它们是大气上层产生的粒子，静止寿命只有约 2.2 微秒，全靠相对论的时间延缓才能活着到达地面（第 38 章）。但真正让本章头疼的还不是它们活了多久，而是它们随时会衰变成别的粒子：在你眼前，粒子的种类和数目一直在变。一套只能描述“固定几个粒子”的理论，连这份参与名单都抄不全。

\step{旧模型遇到麻烦}

现在我们像 1920 年代的物理学家一样，认真地给薛定谔方程“换零件”，看看会在哪里卡住。

第一层麻烦：时间和空间的地位不平等。翻开薛定谔方程（第 44 章），时间在里面只出现一次——波函数对时间的一阶变化率；空间却出现两次——波函数在空间中弯曲程度的二阶导数。时间被放在了特殊的位置上。而第 39 章整章都在讲：相对论里时间和空间必须平权，从一个参考系换到另一个参考系，两者会互相混合。一个把时间和空间区别对待的方程，换一套运动的尺子去量，立刻就变了形。这就像一张地图，南北方向用米作单位、东西方向用尺作单位：你一个人用没问题，一旦要和别人对图，处处都要换算，而且换算规则稀奇古怪。这不是计算麻烦，是结构缺陷。

第二层麻烦：硬换之后，概率变成了负数。那就动手换：把相对论的能量动量关系 \(E^2=p^2c^2+m^2c^4\) 翻译成算符语言（翻译细节见下面的数学桥层）。得到的方程叫克莱因—戈尔登方程，它对时间和空间都是二阶，洛伦兹对称的毛病倒是治好了，却带来两个新病。其一，开平方让能量天然地长出两个分支：\(E=+\sqrt{p^2c^2+m^2c^4}\) 和 \(E=-\sqrt{p^2c^2+m^2c^4}\)。负的能量是什么意思？一个粒子的能量比“什么都没有”还低，那它为什么不一路掉下去、掉到负无穷，顺便释放无穷多的能量？其二更致命：用这个方程去算“在某处找到粒子的概率”，算出来的量在某些区域是负的。概率是负数，这句话根本没有意义——你在抽奖箱里摸到“负三张奖券”是什么体验？注意，这不是哪一步算错了，而是这套翻译本身的问题。顺带一提，这个方程后来被证明并非废物：它是描述自旋为零的粒子的正确方程，只是“职位”不是单粒子波函数，而是量子场。一个方程的失败，有时只是因为它被安在了错误的岗位上。

第三层麻烦，也是最根本的：前两层的修补方案，都悄悄预设了“我们始终在描述同一个粒子”。但相对论一开场就说了 \(E=mc^2\)：能量和质量是同一种东西的两种记账方式，可以互相兑换。只要一次碰撞的能量超过 \(2mc^2\)，两个粒子相撞就可以凭空多出一对粒子来。这不是理论狂想，而是实验事实：加速器里每天都在发生；连两个高能光子相撞都能变成一对电子和正电子（物理学家称之为布赖特—惠勒过程，2021 年由 STAR 实验组用接近光速的金原子核周围的强光场直接观测到）。反过来，一对粒子也可以湮灭成光。一个从头到尾只会说“有一个粒子、永远有一个粒子”的理论，面对这样随时改写参与名单的世界，连剧本的第一页都写不下去。

三个麻烦叠在一起，指向同一个结论：相对论性量子力学不是给旧方程换个零件就能交差的装修工程，而是要把地基挖开重浇。

\step{发明新概念}

1928 年，保罗·狄拉克决定反过来做：不改薛定谔方程的骨架，而是追问——什么样的方程才能既保持“对时间只有一阶”（只有这样，下一时刻的态才完全由现在的态决定，概率才能守恒而且永远非负），又满足相对论的要求？

他把条件写下来，发现这是一个看上去无解的要求：一阶的时间方程，平方之后却必须等于相对论的 \(E^2=p^2c^2+m^2c^4\)。普通数字做不到这一点——没有任何普通系数能让 \((ap+b)^2\) 等于 \(p^2c^2+m^2c^4\) 而不留下难看的交叉项。狄拉克的惊人之举是：既然普通数不行，那就让系数是矩阵。于是波函数不再是一个数，而是一列四个数（后来称为“旋量”），方程里的系数是四个遵循特殊乘法规则的 \(4\times4\) 矩阵。这不是为了炫技，而是被两个原理的夹缝硬生生逼出来的——数学桥层里我们会看到这个逻辑有多紧凑。

方程写出来之后，开始自己“说话”。它给出的东西远远超过了狄拉克下订单时想要的，一共三份意外的礼物。

第一份礼物：自旋自动出现了。解这个方程，粒子天然携带大小为 \(\hbar/2\) 的内禀角动量，不需要像第 47 章那样作为额外的实验事实“贴”进理论里；电子磁矩的大小（\(g\) 因子约等于 2）也自动落出来。斯特恩—盖拉赫实验那个悬了六年的谜——银原子束为什么劈成两束——原来一直藏在这个方程的矩阵结构里。一个为别的目的发明的理论，免费解释了一桩旧案，这在物理学史上是理论“走对了”的强烈信号。

第二份礼物：负能量解还在，但这次狄拉克没有扔掉它。他的处理方案大胆得近乎疯狂（请分清层次：以下是历史上的解释图像，不是实验事实）：他提出，真空中所有负能态早已被电子填满，形成一片看不见的“狄拉克海”；由于泡利不相容原理，普通电子掉不进去，所以世界稳定。而海里如果被敲出一个空穴，这个空穴在各方面都表现得像一个带正电、质量与电子严格相等的粒子。直觉在这里狠狠摔了一跤：狄拉克最初猜这个空穴就是质子——当时已知的唯一带正电的粒子。但空穴的质量必须和电子一模一样，质子却重了 1836 倍，账目怎么都对不上。理论被迫预言了一种全新粒子。1932 年，安德森在宇宙线云室里真的拍到了它：一条径迹，弯曲方式表明它带正电，弯曲程度却表明它和电子一样轻。正电子就这样被找到了。今天我们知道，“海”的图像只是一个历史脚手架，现代理论不再需要它；但“正电子存在”是板上钉钉的实验事实。

第三份礼物：守恒律的重新洗牌。旧世界里“粒子数守恒”退场了——粒子和反粒子可以成对产生、成对湮灭。但更深的守恒律留了下来，而且面目更清晰：能量守恒、动量守恒、电荷守恒。空穴带着与电子相反的电荷，所以产生一对，总电荷仍是零；湮灭一对，总电荷仍是零。守恒的不是“东西的数量”，而是更抽象的账目。

到这里，新概念的轮廓已经浮出来了：让量子力学和相对论真正兼容的代价，是放弃“粒子是最基本的对象”这个信念本身。
